\section{Stage 5: Can the Instrument Be Made Cheaper?}
\label{sec:precision}
\FigurePrecision

\step{Design and predictions.}
The ten seeds are three quarters of the instrument's GPU time.
Unlike distillation (\S\ref{sec:cascade}), reduced precision changes only the arithmetic, so we ran the frozen M1 at reduced precision and compared each format with 16-bit on the same GPU type.
FP8 (8-bit weights and activations) was compared through the full instrument, NF4 (4-bit weights, as in QLoRA, \citealp{dettmers2023qlora}) at the level of M1.
We predicted that FP8 would keep $\ge$98\% label agreement and \Stwo{} within 0.01 (H26-1) and be at least 1.3$\times$ faster at the same batch size (H26-2).
The non-obvious prediction, H26-4, was that FP8's flips would fall where the 27B re-decides, so the instrument would lose less than M1.

\step{Result.}
Both FP8 predictions about M1 failed (Figure~\ref{fig:precision}c): FP8 changes 3.9\% of the ensemble's argmax labels (17 of M1's final decisions) and lowers M1's \Stwo{} by 0.045 with the parameters frozen.
We ran FP8 on the native FP8 tensor cores of an RTX PRO 6000 (Blackwell). As implemented it still lowers batched throughput to 0.83$\times$, because activations are quantized in unfused operations around each FP8 matrix product. Reported FP8 speed-ups rely on kernels that fuse these steps into the matrix product, which we did not use. FP8 also barely changes peak memory (0.96$\times$), since the embeddings, norms and head stay in 16 bit.
Applied to the accuracy point's one-GPU cost, the slower M1 would raise it from 0.538 to 0.619\,\mbox{GPU-s} per item (+15\%).
H26-4 held on its two registered counts (Figure~\ref{fig:precision}a,b): 71\% of FP8's 17 flips land on items the 27B re-decides, twice the chance rate, and the system flips 16 decisions against M1's 17.
All 10 of M1's membership-changing flips lose reference-set membership, whereas the system's go both ways (5 lose, 7 gain), so with the same frozen parameters it loses 0.003 against M1's 0.045.
The cascade does not undo the flips but turns a one-sided loss into noise.
NF4 changes more (86.7\% ensemble agreement) and meets our registered QLoRA trigger.
It saves memory but not time: peak memory falls to 0.67$\times$ (16.8 to 11.3\,GB), short of the 0.4$\times$ we predicted, while batch-1 latency rises from 169 to 259\,ms, because the 4-bit weights are stored compactly but computed in bf16.
Its throughput (0.82$\times$) is not like for like, since the 16-bit reference ran split over two cards.
Its cascade comparison is incomplete (27B records are missing for 60 of 136 deferred items), so we draw no conclusion from it.
Precision is not a cost lever with these kernels.
